
\section{Samples}
\label{sec:manual-samples}

This chapter contains documentation about the sample programs provided with Cubed\-OS. These
programs illustrate various aspects of using CubedOS in the real world that might not be
covered by the more formal documentation presented so far. New users of CubedOS are encouraged to
study, build, execute, and modify the sample programs as a way of learning about the system and
how it can be used.

\subsection{Echo Sample}
\label{sec:manual-echo-sample}

This section describes the Echo Sample that can be found in the \filename{samples/echo} folder.
This sample exemplifies a simple message/reply demo. These messages take place as simple pings
back and forth between a Client module and a Server module.

The main.adb of this program has several with clauses to include |Echo_Client.Messages| and
|Echo_Server.Messages|. It also uses |Log_Server.Messages| in order for the modules to be
capable of logging various errors or informational output. This main file contains a single
procedure of first priority with a single loop that has a small delay. This loop spends most of
its time sleeping, and the real work is done inside the Client and Server modules which are
unreferenced in the main itself.

Both |Echo_Client.ads| and |Echo_Server.ads| have with clauses for |Message_Manager|. These are
of course unreferenced, but are very important. The package |Message_Manager| creates a new
instance with a domain number, module count, mailbox size, and maximum message size. This sets
the environment for all future message sending. This file also elaborates
|CubedOS.Generic_Message_Manager|, which contains all necessary functions and procedures for
managing messages.

|Echo_Client.Messages.adb| implements the main part of its module. At its core, this component
sends an initial request message to the server, and continues to send messages each time it gets
a reply back. There is a single task called |Message_Loop| in the body. This task first runs the
initialize procedure, and then loops while waiting for messages. Inside the initialize procedure
is where the first message gets sent. This procedure uses the |Echo_Server.API| to encode a ping
request message. The sender address is set to |Name_Resolver.Echo_Client|. The
|Name_Resolver.ads| file contains the message addresses for the Client, Server, and Log Server.
After initializing, the task starts looping. In each loop, it waits for a message. Once it
receives a message, it processes that message and sends another.

|Echo_Server.Messages.adb| Has a similar |Message_Loop|. This task does not send an initial
message, but instead, has just a loop. The loop waits for a message right away. When one is
received, it processes it and sends a reply message using |Echo_Server.API|. Then, it starts the
Message loop over again.

This goes on infinitely. The log server logs information every time a reply is received by the
Client. It also logs errors when a message goes wrong. This is to show the user the basic
functionality of the log server, and how it can be used with the message manager.


\subsection{LineRider Sample}
\label{sec:manual-linerider-sample}

\todo{Finish Me!}


\subsection{Pathfinder Sample}
\label{sec:manual-pathfinder-sample}

This sample is intended to demonstrate the problem of priority inversion in CubedOS. It is
loosely based on Mars Pathfinder, which had an actual issue with priority inversion during the
mission.

Section~\ref{sec:architecture-message-priorities} describes the priority architecture of
CubedOS. Briefly, every module has a statically assigned priority based on the priority given to
the message loop that fetches and processes messages. There is currently no facility to change
module (task) priorities once they are assigned. Priority inversion can happen in the following
circumstances:

\begin{itemize}

\item A low priority module begins processing a message.

\item A medium priority module becomes ready to run, and the runtime system suspends the low
priority module in favor of the medium priority one.

\item A high priority module becomes ready to run, and the runtime system suspends the medium
priority module in favor of the high priority one.

\item The high priority module sends a message to the low priority module. This message requires
a reply, so the high priority module waits for that reply to arrive by blocking on fetching from
its mailbox.

\item The runtime system resumes the medium priority module which, unfortunately in this
hypothetical case, runs for an indefinitely long time.

\end{itemize}

In this scenario, the low priority module is never given a chance to finish what it started,
never sees the message from the high priority module, and never replies to that message. The
high priority module is now waiting, in effect, on the medium priority module to finish, since
that is what is holding up the low priority module from making progress. Yet high priority tasks
should not be held up by the execution of any task with lower priority. This is priority
inversion.

\pet{Finish this! More needs to be said about the sample and how to interpret its output. It
would also be good to include some citation(s) regarding the actual Pathfinder issue.}


\subsection{STM32F4 Sample}
\label{sec:manual-stm32f4-sample}

This section describes the architecture of the STM32F4 sample that can be found in a folder of
the same name in the CubedOS repository. This sample makes use of a simple evaluation board from
STMicroelectronics. The README file in the STM32F4 folder gives specifics about the board and
about how to set up the development environment for it.

The sample program is simple in effect, but written in a way as to illustrate some of CubedOS's
structure and features. The program flashes the LEDs on the board in a rotating sequence, from
green to orange to red to blue. The program also watches the user button on the board (the blue
button). When pressed, the speed of the flashing increases in steps, from low to medium to high.
The program has no other features.

This program illustrates several things.
\begin{itemize}

  \item Using CubedOS with a non-native CPU. The sample was developed on a Linux desktop system
  running a 64~bit Intel processor. The project, however, is configured to generate code for an
  ARM Cortex M4 CPU. This affects not only the compiler being used, but also the \SPARK\
  configuration used to analyze the code.

  \item Using a small, bare board embedded platform. The sample makes use of AdaCore's light
  tasking runtime for the STM32F4 along with a few files provided by AdaCore for accessing the
  hardware on the board. It does not make use of any conventional operating system.

  \item Using CubedOS modules as drivers for onboard hardware. Although this simple application
  could be more easily programmed without creating driver modules, for purposes of illustration
  it includes two driver modules for the hardware used: an LED driver and a button driver.

  \item Using a CubedOS ``controller'' module as the main module of the program. The controller
  module responds to events generated on the board by sending messages to other modules, as
  appropriate. The controller serves as the main part of the application and encodes the high
  level application logic. The Ada main procedure contains no application functionality at all.
  This is a typical for CubedOS applications.

  \item Using the publish-subscribe server, a core CubedOS module, for decoupling event sources
  (as generated by the button driver) to event sinks (the controller module). This allows the
  driver modules to be reusable in other applications for this platform. Specifically the button
  driver does not directly interact with the application-specific controller. It could be
  dropped into another application, along with the core publish-subscribe server module, without
  modification.

\end{itemize}

In the subsections that follow more details are given about each of these points.

\subsubsection{Non-Native CPU}

The STM32F4 sample uses an ARM Cortex M4 processor. This entails installing the necessary Ada
cross-compiler from AdaCore. The project file specifies the tool chain to be used. The other
important aspect of using a non-native CPU is configuring \SPARK\ to know about the target
platform attributes. Since the sizes and ranges of the primitive types are often different on
the embedded platform than on the native platform, the sizes and ranges of the Ada base types
are also often different. This affects when computations might overflow. Consider, for example,
this simple case:

\begin{lstlisting}
declare
   type Some_Type is range 0 .. 20_000;
   A, B, C : Some_Type;
begin
   -- etc...
   A := (B + C) / 2;
end;
\end{lstlisting}

On a system where the base type of |Some_Type| has 32~bits, there is no possibility of overflow
when computing |(B + C)/2| since, at most, |B + C| would be 40,000. However, if a 16~bit base
type is used, as might be the case on a small microcontroller platform, |B + C| might go outside
the range of 16-bit signed integers, and an overflow would be possible. The \SPARK\ tools know
what base types are used by the Ada compiler, but that depends on the platform being targeted
by that compiler. It is thus essential to inform the \SPARK\ tools about the target platform so
that the proofs are relevant and correct for the target context. By default, the \SPARK\ tools
will assume the native platform's properties, which might not be appropriate when cross
compiling.

\todo{Finish Me!}
